\documentclass[a4paper,10pt]{article}
\usepackage[utf8]{inputenc}
\usepackage[margin=1.5cm, bottom=2cm]{geometry}
\usepackage{fancyhdr}
\usepackage[table]{xcolor}
\usepackage{enumitem}
\usepackage{xcolor}
\usepackage{makeidx}
\usepackage[most]{tcolorbox}
\newtcolorbox{osservazioni}{
  colback=gray!8,    % sfondo grigio chiaro
  colframe=white,      % nessun bordo
  arc=2mm,            % angoli smussati
  boxrule=0pt,        % spessore bordo 0
  fontupper=\footnotesize,
  left=2mm, right=2mm, top=2mm, bottom=2mm,
  before skip=7pt,   % niente spazio prima
  after skip=15pt     % niente spazio dopo
}
\newtcolorbox{osservazioni2}{
  colback=gray!8,    % sfondo grigio chiaro
  colframe=white,      % nessun bordo
  arc=2mm,            % angoli smussati
  boxrule=0pt,        % spessore bordo 0
  fontupper=\small,
  left=5mm, right=5mm, top=3mm, bottom=3mm,
  before skip=0pt,   % niente spazio prima
  after skip=30pt     % niente spazio dopo
}
\pagestyle{fancy}
\usepackage{tikz}
\usetikzlibrary{automata, positioning}
\pagestyle{fancy}

\usepackage{makeidx}

%%%%%%%%%%%%%%%%%%%
\newcommand{\EsercitazioneNum}{09}
\newcommand{\Data}{04/11/2025}
\newcommand{\DocType}{Esercitazione} % o "Eserciziario"
\newcommand{\FullTitle}{\DocType\ \EsercitazioneNum}

\newcommand{\Header}{
\noindent
  \small \textbf{Computabilità, Complessità e Logica - a.a. 2025/26} \hfill \textbf{\FullTitle} \\[0.2cm]
  \scriptsize Matteo Liotta, \texttt{matteo.liotta@studenti.units.it} \hfill \Data \\[0.2cm]
  \noindent\makebox[\linewidth]{\rule{\linewidth}{0.4pt}} \\[0.3cm]}

\newcommand{\NewPageHeader}{
  \vfill
  \noindent\makebox[\linewidth]{\rule{\linewidth}{0.4pt}}
  \newpage
  {\Header}
}
%%%%%%%%%%%%%%%%%%%

% Numero di pagina in basso al centro
\fancyhf{}
\fancyfoot[C] \thepage
\renewcommand{\footrulewidth}{0pt} % niente linea in basso
\renewcommand{\headrulewidth}{0pt} % niente linea in alto

\begin{document}

\footnotesize  % font generale più piccolo
\Header\\[0.7cm]


%%%%%%%%%%%%%%%% FOGLIO 2


%% INDEX
\addcontentsline{toc}{section}{Esercitazione 07: \textit{Macchine di Turing}}

\noindent\Large \textbf{Esercitazione 9}: \textit{Macchine di Turing (cont'd)}
\begin{osservazioni}
    \textbf{\textit{Nota}}. Sono indicati in \textcolor{blue}{\textit{blu}} gli esercizi che non saranno trattati durante l'esercitazione, bensì lasciati come strumento aggiuntivo di esercitazione individuale. Resta comunque assoluta disponibilità per la loro risoluzione negli incontri successivi, qualora necessario.
\end{osservazioni}

{\normalsize
\begin{itemize}[label=, leftmargin=0pt]

    %% INDEX
    \addcontentsline{toc}{subsection}{Esercizio 52.} 
    %%
    \item \textbf{Esercizio 52 (*)}\\ Sia $L$ linguaggio su $\Sigma=\{a,b\}$ (cfr. \textit{approfondimento 02}), tale per cui:
    {\begin{center}
        $L=\{w\#w'| w,w'\in\Sigma^*, \quad \exists i \in \mathbf{N} \quad t.c.\quad  w(i)=w'(i)\ \}$
    \end{center}} 
    Si rappresenti la macchina di Turing tale da riconoscere il linguaggio $L$.\\[0.1cm]

    
    %% INDEX
    \addcontentsline{toc}{subsection}{Esercizio 53.} 
    %%
    \item \textbf{\textcolor{black}{Esercizio 53.} (*)}\\ Sia $L$ linguaggio su $\Sigma=\{a,b,c\}$, tale per cui 
    {\begin{center}
        $L = \{ x\#y \mid x,y\in\Sigma^*\quad e \quad x \text{ prefisso di }y\}$
    \end{center}}
    Si rappresenti la macchina di Turing tale da riconoscere il linguaggio $L$.\\[0.1cm]

    
    %% INDEX
    \addcontentsline{toc}{subsection}{Esercizio 54.} 
    %%
    \item \textbf{\textcolor{blue}{Esercizio 54.} (*)}\\ Sia $L$ linguaggio su $\Sigma=\{a,b,c\}$, tale per cui 
    {\begin{center}
        $L = \{ x\#y \mid x,y\in\Sigma^*\quad e \quad x \text{ sottostringa di }y\}$
    \end{center}}
    Si rappresenti la macchina di Turing tale da riconoscere il linguaggio $L$.\\[0.1cm]
    
    %% INDEX
    \addcontentsline{toc}{subsection}{Esercizio 55.} 
    %%
    \item \textbf{\textcolor{orange}{Esercizio 55.} (***)}\\ 
    Si considerino le seguenti identificazioni, considerando il seguente alfabeto $\{0,1\}$ e simboli ausiliari $\{\$_1,\$_2,\$_c,\#_1,\#_2,\sqcup\}$:
    \begin{itemize}[label=\textbullet]
        \item \texttt{00} $\equiv$ \texttt{IF} $\quad \mid \quad $ \texttt{01} $\equiv$ \texttt{ELSE}
        \item \texttt{0} $\equiv$ \texttt{$\leq$} $\quad \mid \quad $ \texttt{1} $\equiv$ \texttt{$>$}
        \item \texttt{10} $\equiv$ \texttt{AND} $\quad \mid \quad $ \texttt{1} $\equiv$ \texttt{OR (element-wise)}
    \end{itemize}
    Dove sono presenti le seguenti regole sintattiche:
    \begin{itemize}[label=\textbullet]
        \item Ogni operazione di controllo con if richiede l'indicazione del termine di paragone voluto, i termini da confrontare, correttamente separati:\\ \texttt{IF [binary num 1]  [binary num 2] CONFRONTO THEN [what to do]} $\equiv$ 
        \texttt{IF$\$_1x\$_2y\$_cc\#_1command\#_2$} \\{$x,y\in\Sigma^*$}
        \item Si definisce l'operazione alternativa, posta dopo la codifica del THEN dell'IF:\\ \texttt{ELSE [what to do]} $\equiv$ \texttt{ELSE$\#_1command\#_2$}
        \item Mentre, con operazioni (then), in questo caso si sono considerate somme e sottrazioni elementwise degi termini:\\ \texttt{OPERATION [binary num 1] [binary num 2]} $\equiv$ 
        \texttt{OP$\$_1x\$_2y\#_2 \quad x,y\in\Sigma^*$}
        
    \end{itemize}
    Si scriva una macchina di Turing che, ricevendo in input una stringa del programma (vedi esempio), ritorni sul nastro il risultato finale dell'(ultima) operazione

\item \textit{e.g.} \textit{Se 010 $>$ 001 allora effettuane l'AND}: \begin{itemize} \item \texttt{INPUT: 00$\$_1$010$\$_2$001$\$_c$1$\#_1$10$\$_1$010$\$_2$001$\#_2$} \item \texttt{OUTPUT: 011} \end{itemize}







\vfill
\noindent\makebox[\linewidth]{\rule{\linewidth}{0.1pt}}



\end{document}